\section{Integration with TGCE Loss}

The model produces two distinct outputs: segmentation masks
$\mathbf{\hat{Y}} = f(\mathbf{X};\theta)$ and reliability maps
$\{\Lambda_r(\mathbf{X};\theta)\}_{r=1}^R$. These outputs work in
tandem with the $TGCE_{SS}$ loss function described in Section
\ref{sec:proposed-loss-functions}. The loss function simultaneously guides the
learning of both the segmentation masks and reliability maps,
ensuring that the model learns to balance the contributions of
different annotators based on their estimated reliability.

The TGCE$_{SS}$ loss function is implemented differently for each reliability
approach, with specific adaptations to handle the different spatial dimensions
of the reliability maps. Naturally, the model itself container tunable
hyperparameters related to the \gls{GCE} - \gls{MAE} trade off behavior
(also known as $q$ parameter) and the unreliability tolerance $\tilde C$,
however, some additional hyperpamarameters were added with the purpose of
regularizing the overall loss for every experiment.

\subsection{Regularization terms}\label{sec:regularizers}

All three implementations include regularization terms to ensure
proper learning of the reliability maps:

\subsubsection{Numerical stability}
To ensure numerical stability in the loss function computations,
particularly in division operations, a small epsilon value $\epsilon$
is added to denominators. This prevents division by zero
errors and stabilizes gradient calculations during training. The
epsilon value is kept small enough to not significantly impact the
loss function behavior while providing the necessary numerical safeguards.

\subsubsection{Penalty for extreme reliabilities}

Some optimization processes might favor a convergence to extreme reliability
estimations, by assigning either 0 or 1 to all labelers. To avoid this, a
regularization term is added to the loss function which penalizes the model
for assigning extreme reliability values and favoring those around a
centroid. High scaling values for this regularization term might
lead to a convergence to a very homogeneous reliability map across all
labelers, which is not desired either.

\begin{equation}
  \mathcal{L}_{centroid} = \alpha_{centroid} \cdot
  \mathbb{E}[(\Lambda_r - 0.5)^2]
\end{equation}

\subsubsection{Entropy regularization}
In this case, entropy is useful to ensure that the reliability maps
are not too homogeneous. The entropy regularization term encourages
the model to assign more discriminative reliability values by penalizing
uniform distributions. This prevents the model from converging to
uninformative reliability estimates where all labelers appear equally
reliable, thus maintaining meaningful distinctions between more and
less trustworthy annotators.

\begin{equation}
  \mathcal{L}_{entropy} = \alpha_{entropy} \cdot \mathbb{E}[\Lambda_r
    \log(1 + \Lambda_r)
  +(1 - \Lambda_r) \log(1 + (1 - \Lambda_r))]
\end{equation}

\subsubsection{Sum-to-one regularization}

The reliability branch for scalar, pixel and features approaches,
might or might not
be forced to sum to one through a softmax operation \footnote{This
  was left as flexible and configurable in the Network Architecture
Building API from seg\_tgce}. For the cases in which this sum is not
forced to be unitary the flexibility of this sum is manageable
through a regularization term which penalizes the model for not
summing to one. This was found useful in situations in which the
reliability maps were not expected to be punctuated (e.g. for the
scalar reliability) as $1/R$, but rather close to the unit for all of
them (if all labelers were good). This regularization term was found
particularly useful in situations where labelers have good and
similar levels of expertise.

\begin{equation}
  \mathcal{L}_{sum} = \alpha_{sum} \cdot
  \mathbb{E}[\sum_{r\in R}( \Lambda_r - 1)^2]
\end{equation}

\subsubsection{Missing labels penalizing}
In situations were the datasets has a low labeling rate, the model might
take an optimization shortcut by assigning high reliability values to
labelers which did not label at all, since the element-wise product with the
mask would be close to zero, which favors a minimization target,
which is not desired. To avoid this, a term was added to the loss
function which penalizes the model for assigning high reliability
values to unlabeled patches, as follows:

\begin{equation}
  \mathcal{L}_{\lambda} = \gamma \max \left(0, \sum_{\forall r \in R}
    \frac{1}{1-\alpha_{\lambda}} \left(\Lambda_r - \mu_r\right) \cdot
  e^{\frac{\Lambda_r - \mu_r - 1}{\beta}}\right)
\end{equation}

Where $\Lambda_r$ is the estimated reliability of the labeler $r$ for a
certain patch, $\mu_r$ is the one hot label mask of the labeler
$r$ for the same patch, $\alpha_{\lambda}$ and $\beta$ are hyperparameters
which control the penalty strength for high and low estimated
reliabilities, respectively and $\gamma$ is a hyperparameter
which controls the overall strength of the penalty.

\begin{figure}[H]
  \begin{tikzpicture}
    \begin{axis}[
        axis lines=middle,
        xlabel={$z=\Lambda_r - \mu_r$},
        ylabel={$\mathcal{L}_{\lambda}(z)$},
        width=10cm, height=7cm,
        domain=-2:4,
        samples=400,
        smooth,
        thick,
        ymin=0, ymax=5,         % <-- cap y range
        xtick={-2,-1,0,1,2,3,4},
        ytick={0,1,2,3,4,5},
        every axis x label/.style={at={(ticklabel* cs:1.05)}, anchor=west},
        every axis y label/.style={at={(ticklabel* cs:1.05)}, anchor=south},
      ]

      % Parameters
      \def\alphaval{0.2}
      \def\betaval{1.0}   % <- larger beta keeps growth moderate

      % Plot: max(0, expression)
      \addplot[blue,thick]
      {max(0, (1/(1-\alphaval))*(x)*exp((x - 1)/\betaval))};

    \end{axis}
  \end{tikzpicture}
  \caption{Example of the $\mathcal{L}_{\lambda}(z)$ function for
  values of $\alpha_{\lambda}=0.2$ and $\beta=1.0$.}
  \label{fig:missing_labels_penalizing}
\end{figure}

This function applies no penalty for cases in which  $\Lambda_r <
\mu_r$, meaning no penalty for assigning any reliability value
(ranging from 0 to 1) to a labeler who actually labeled a given patch.
On the other hand, a penalty is indeed applied in cases where
$\Lambda_r > \mu_r$, meaning a penalty for assigning high reliability
values to a labeler which did not label at all. The proportion on how
high and how quickly (in the space of difference $\Lambda_r - \mu_r$)
the penalty is applied is controlled by the hyperparameters $\alpha_{\lambda}$
and $\beta$. The behavior of this term is shown in Figure
\ref{fig:missing_labels_penalizing}.

These regularization terms lead to the overall loss function being in the form:

\begin{align}
  L_{TGCEreg} &= L_{TGCEss} \\
  &+\alpha_{reg} \cdot \mathbb{E}[(\Lambda_r - 0.5)^2] \nonumber \\
  &+\alpha_{entropy} \cdot \mathbb{E}[\Lambda_r \log(1 + \Lambda_r)
  +(1 - \Lambda_r) \log(1 + (1 - \Lambda_r))] \nonumber \\
  &+\alpha_{sum} \cdot \mathbb{E}[\sum_{\forall r \in R}(\Lambda_r -
  1)^2] \nonumber \\
  &+\gamma \cdot \max \left(0, \sum_{\forall r \in R}
    \frac{1}{1-\alpha_{\lambda}} \left(\Lambda_r - \mu_r\right) \cdot
  e^{\frac{\Lambda_r - \mu_r - 1}{\beta}}\right) \nonumber
\end{align}

For the scalar reliability approach, the loss function operates on
global reliability values per annotator. The reliability maps are
expanded to match the spatial dimensions of the predictions through
tiling operations. This approach is computationally efficient but
provides the same reliability value across all spatial locations.

The feature-based reliability approach requires additional processing
to handle the spatial dimensions of the bottleneck features. The
reliability maps are resized using bilinear interpolation to match
the spatial dimensions of the predictions. This approach provides
coarse-grained spatial reliability information while maintaining
computational efficiency.

The pixel-wise reliability approach operates directly on full-resolution
reliability maps, requiring no spatial transformations. This approach
provides the most detailed spatial reliability information but requires
more computational resources.
